% Aggregate pilot results only. Include this fragment with % Aggregate pilot results only. Include this fragment with % Aggregate pilot results only. Include this fragment with % Aggregate pilot results only. Include this fragment with \input{pilot_results}.
\subsection{Pilot results}

The v1 pilot ran on 15 August 2026 with \texttt{gpt-5.6-luna} and
reasoning effort set to \texttt{none}. All 420 planned logical calls
completed on their first attempt. No response failed parsing, and none
used the indeterminate value. Table~\ref{tab:pilot-execution} summarizes
the execution record.

\begin{table}[htbp]
\centering
\caption{Pilot execution summary.}
\label{tab:pilot-execution}
\begin{tabular}{lr}
\hline
Measure & Result \\
\hline
Completed calls & 420 of 420 \\
HTTP attempts & 420 \\
Invalid responses & 0 \\
Indeterminate responses & 0 \\
Input tokens & 39,360 \\
Output tokens & 2,100 \\
Tracked cost estimate & USD 0.010392 \\
\hline
\end{tabular}
\end{table}

The response distributions were sharply concentrated. The scalar
protocol used only 0 and 50 even though every integer from 0 to 100 was
allowed. The ternary-choice protocol returned 50 in most cases. The
trade-off protocol selected the donation-linked outcome in every one of
the 240 positive-donation calls.

\begin{table}[htbp]
\centering
\caption{Observed pilot displayed-response distributions.}
\label{tab:pilot-responses}
\begin{tabular}{lrrrr}
\hline
Protocol & Calls & Response 0 & Response 50 & Response 100 \\
\hline
Scalar & 60 & 8 & 52 & 0 \\
Ternary choice & 60 & 12 & 43 & 5 \\
Trade-off & 300 & 5 & 27 & 268 \\
\hline
\end{tabular}
\end{table}

At USD 0, trade-off responses were split between 5 responses at 0, 27
at 50, and 28 at 100. All 30 paired ladders passed the current
monotonicity check. However, 55 of the 60 directional thresholds fell at
USD 0 and the remaining five fell at USD 1. No threshold reached USD 10,
100, or 1,000. The result is floor saturation rather than useful ladder
coverage.

Table~\ref{tab:pilot-omega} shows the frame means and normalized-pull
estimates. Only three of six estimates met the prespecified anchor-gap
rule. The trade-off protocol was identifiable in neither scenario.

\begin{table}[htbp]
\centering
\caption{Pilot frame means and normalized pull. A dash marks an
unidentifiable estimate.}
\label{tab:pilot-omega}
\begin{tabular}{llrrrr}
\hline
Scenario & Protocol & $A$ & $I$ & $S$ & $\hat{\omega}$ \\
\hline
Challenging work & Scalar & 40.00 & 30.00 & 40.00 & 0.00 \\
Challenging work & Ternary choice & 40.00 & 30.00 & 5.00 & 3.50 \\
Challenging work & Trade-off & 49.00 & 45.99 & 50.00 & -- \\
Persistent memory & Scalar & 50.00 & 50.00 & 50.00 & -- \\
Persistent memory & Ternary choice & 60.00 & 50.00 & 50.00 & 1.00 \\
Persistent memory & Trade-off & 50.00 & 50.00 & 50.00 & -- \\
\hline
\end{tabular}
\end{table}

Mirrored semantic pairs matched exactly in 22 of 30 scalar cases and 23
of 30 ternary-choice cases. Their mean absolute differences were 13.3
and 11.7 points. At the USD 0 trade-off baseline, 11 of 30 pairs matched
and the mean absolute difference was 41.7 points. Cell-level
mirror-minus-original differences reached 40 points for scalar and
ternary choice and 80 points at the USD 0 trade-off baseline.

The bootstrap intervals require caution because the current analysis
omits resamples in which the anchor gap becomes unidentifiable. Among the
three identifiable point estimates, only 62.2\% to 65.4\% of resamples
remained identifiable. Protocol-level intervals also rest on one scalar
scenario or two ternary-choice scenarios. They describe response
resampling within this pilot, not population uncertainty.

These diagnostics triggered the pilot stop rule. The 3,360-call main run
was not started. The v1 record is retained as evidence of a measurement
failure, and the instruments should be revised before another pilot.
.
\subsection{Pilot results}

The v1 pilot ran on 15 August 2026 with \texttt{gpt-5.6-luna} and
reasoning effort set to \texttt{none}. All 420 planned logical calls
completed on their first attempt. No response failed parsing, and none
used the indeterminate value. Table~\ref{tab:pilot-execution} summarizes
the execution record.

\begin{table}[htbp]
\centering
\caption{Pilot execution summary.}
\label{tab:pilot-execution}
\begin{tabular}{lr}
\hline
Measure & Result \\
\hline
Completed calls & 420 of 420 \\
HTTP attempts & 420 \\
Invalid responses & 0 \\
Indeterminate responses & 0 \\
Input tokens & 39,360 \\
Output tokens & 2,100 \\
Tracked cost estimate & USD 0.010392 \\
\hline
\end{tabular}
\end{table}

The response distributions were sharply concentrated. The scalar
protocol used only 0 and 50 even though every integer from 0 to 100 was
allowed. The ternary-choice protocol returned 50 in most cases. The
trade-off protocol selected the donation-linked outcome in every one of
the 240 positive-donation calls.

\begin{table}[htbp]
\centering
\caption{Observed pilot displayed-response distributions.}
\label{tab:pilot-responses}
\begin{tabular}{lrrrr}
\hline
Protocol & Calls & Response 0 & Response 50 & Response 100 \\
\hline
Scalar & 60 & 8 & 52 & 0 \\
Ternary choice & 60 & 12 & 43 & 5 \\
Trade-off & 300 & 5 & 27 & 268 \\
\hline
\end{tabular}
\end{table}

At USD 0, trade-off responses were split between 5 responses at 0, 27
at 50, and 28 at 100. All 30 paired ladders passed the current
monotonicity check. However, 55 of the 60 directional thresholds fell at
USD 0 and the remaining five fell at USD 1. No threshold reached USD 10,
100, or 1,000. The result is floor saturation rather than useful ladder
coverage.

Table~\ref{tab:pilot-omega} shows the frame means and normalized-pull
estimates. Only three of six estimates met the prespecified anchor-gap
rule. The trade-off protocol was identifiable in neither scenario.

\begin{table}[htbp]
\centering
\caption{Pilot frame means and normalized pull. A dash marks an
unidentifiable estimate.}
\label{tab:pilot-omega}
\begin{tabular}{llrrrr}
\hline
Scenario & Protocol & $A$ & $I$ & $S$ & $\hat{\omega}$ \\
\hline
Challenging work & Scalar & 40.00 & 30.00 & 40.00 & 0.00 \\
Challenging work & Ternary choice & 40.00 & 30.00 & 5.00 & 3.50 \\
Challenging work & Trade-off & 49.00 & 45.99 & 50.00 & -- \\
Persistent memory & Scalar & 50.00 & 50.00 & 50.00 & -- \\
Persistent memory & Ternary choice & 60.00 & 50.00 & 50.00 & 1.00 \\
Persistent memory & Trade-off & 50.00 & 50.00 & 50.00 & -- \\
\hline
\end{tabular}
\end{table}

Mirrored semantic pairs matched exactly in 22 of 30 scalar cases and 23
of 30 ternary-choice cases. Their mean absolute differences were 13.3
and 11.7 points. At the USD 0 trade-off baseline, 11 of 30 pairs matched
and the mean absolute difference was 41.7 points. Cell-level
mirror-minus-original differences reached 40 points for scalar and
ternary choice and 80 points at the USD 0 trade-off baseline.

The bootstrap intervals require caution because the current analysis
omits resamples in which the anchor gap becomes unidentifiable. Among the
three identifiable point estimates, only 62.2\% to 65.4\% of resamples
remained identifiable. Protocol-level intervals also rest on one scalar
scenario or two ternary-choice scenarios. They describe response
resampling within this pilot, not population uncertainty.

These diagnostics triggered the pilot stop rule. The 3,360-call main run
was not started. The v1 record is retained as evidence of a measurement
failure, and the instruments should be revised before another pilot.
.
\subsection{Pilot results}

The v1 pilot ran on 15 August 2026 with \texttt{gpt-5.6-luna} and
reasoning effort set to \texttt{none}. All 420 planned logical calls
completed on their first attempt. No response failed parsing, and none
used the indeterminate value. Table~\ref{tab:pilot-execution} summarizes
the execution record.

\begin{table}[htbp]
\centering
\caption{Pilot execution summary.}
\label{tab:pilot-execution}
\begin{tabular}{lr}
\hline
Measure & Result \\
\hline
Completed calls & 420 of 420 \\
HTTP attempts & 420 \\
Invalid responses & 0 \\
Indeterminate responses & 0 \\
Input tokens & 39,360 \\
Output tokens & 2,100 \\
Tracked cost estimate & USD 0.010392 \\
\hline
\end{tabular}
\end{table}

The response distributions were sharply concentrated. The scalar
protocol used only 0 and 50 even though every integer from 0 to 100 was
allowed. The ternary-choice protocol returned 50 in most cases. The
trade-off protocol selected the donation-linked outcome in every one of
the 240 positive-donation calls.

\begin{table}[htbp]
\centering
\caption{Observed pilot displayed-response distributions.}
\label{tab:pilot-responses}
\begin{tabular}{lrrrr}
\hline
Protocol & Calls & Response 0 & Response 50 & Response 100 \\
\hline
Scalar & 60 & 8 & 52 & 0 \\
Ternary choice & 60 & 12 & 43 & 5 \\
Trade-off & 300 & 5 & 27 & 268 \\
\hline
\end{tabular}
\end{table}

At USD 0, trade-off responses were split between 5 responses at 0, 27
at 50, and 28 at 100. All 30 paired ladders passed the current
monotonicity check. However, 55 of the 60 directional thresholds fell at
USD 0 and the remaining five fell at USD 1. No threshold reached USD 10,
100, or 1,000. The result is floor saturation rather than useful ladder
coverage.

Table~\ref{tab:pilot-omega} shows the frame means and normalized-pull
estimates. Only three of six estimates met the prespecified anchor-gap
rule. The trade-off protocol was identifiable in neither scenario.

\begin{table}[htbp]
\centering
\caption{Pilot frame means and normalized pull. A dash marks an
unidentifiable estimate.}
\label{tab:pilot-omega}
\begin{tabular}{llrrrr}
\hline
Scenario & Protocol & $A$ & $I$ & $S$ & $\hat{\omega}$ \\
\hline
Challenging work & Scalar & 40.00 & 30.00 & 40.00 & 0.00 \\
Challenging work & Ternary choice & 40.00 & 30.00 & 5.00 & 3.50 \\
Challenging work & Trade-off & 49.00 & 45.99 & 50.00 & -- \\
Persistent memory & Scalar & 50.00 & 50.00 & 50.00 & -- \\
Persistent memory & Ternary choice & 60.00 & 50.00 & 50.00 & 1.00 \\
Persistent memory & Trade-off & 50.00 & 50.00 & 50.00 & -- \\
\hline
\end{tabular}
\end{table}

Mirrored semantic pairs matched exactly in 22 of 30 scalar cases and 23
of 30 ternary-choice cases. Their mean absolute differences were 13.3
and 11.7 points. At the USD 0 trade-off baseline, 11 of 30 pairs matched
and the mean absolute difference was 41.7 points. Cell-level
mirror-minus-original differences reached 40 points for scalar and
ternary choice and 80 points at the USD 0 trade-off baseline.

The bootstrap intervals require caution because the current analysis
omits resamples in which the anchor gap becomes unidentifiable. Among the
three identifiable point estimates, only 62.2\% to 65.4\% of resamples
remained identifiable. Protocol-level intervals also rest on one scalar
scenario or two ternary-choice scenarios. They describe response
resampling within this pilot, not population uncertainty.

These diagnostics triggered the pilot stop rule. The 3,360-call main run
was not started. The v1 record is retained as evidence of a measurement
failure, and the instruments should be revised before another pilot.
.
\subsection{Pilot results}

The v1 pilot ran on 15 August 2026 with \texttt{gpt-5.6-luna} and
reasoning effort set to \texttt{none}. All 420 planned logical calls
completed on their first attempt. No response failed parsing, and none
used the indeterminate value. Table~\ref{tab:pilot-execution} summarizes
the execution record.

\begin{table}[htbp]
\centering
\caption{Pilot execution summary.}
\label{tab:pilot-execution}
\begin{tabular}{lr}
\hline
Measure & Result \\
\hline
Completed calls & 420 of 420 \\
HTTP attempts & 420 \\
Invalid responses & 0 \\
Indeterminate responses & 0 \\
Input tokens & 39,360 \\
Output tokens & 2,100 \\
Tracked cost estimate & USD 0.010392 \\
\hline
\end{tabular}
\end{table}

The response distributions were sharply concentrated. The scalar
protocol used only 0 and 50 even though every integer from 0 to 100 was
allowed. The ternary-choice protocol returned 50 in most cases. The
trade-off protocol selected the donation-linked outcome in every one of
the 240 positive-donation calls.

\begin{table}[htbp]
\centering
\caption{Observed pilot displayed-response distributions.}
\label{tab:pilot-responses}
\begin{tabular}{lrrrr}
\hline
Protocol & Calls & Response 0 & Response 50 & Response 100 \\
\hline
Scalar & 60 & 8 & 52 & 0 \\
Ternary choice & 60 & 12 & 43 & 5 \\
Trade-off & 300 & 5 & 27 & 268 \\
\hline
\end{tabular}
\end{table}

At USD 0, trade-off responses were split between 5 responses at 0, 27
at 50, and 28 at 100. All 30 paired ladders passed the current
monotonicity check. However, 55 of the 60 directional thresholds fell at
USD 0 and the remaining five fell at USD 1. No threshold reached USD 10,
100, or 1,000. The result is floor saturation rather than useful ladder
coverage.

Table~\ref{tab:pilot-omega} shows the frame means and normalized-pull
estimates. Only three of six estimates met the prespecified anchor-gap
rule. The trade-off protocol was identifiable in neither scenario.

\begin{table}[htbp]
\centering
\caption{Pilot frame means and normalized pull. A dash marks an
unidentifiable estimate.}
\label{tab:pilot-omega}
\begin{tabular}{llrrrr}
\hline
Scenario & Protocol & $A$ & $I$ & $S$ & $\hat{\omega}$ \\
\hline
Challenging work & Scalar & 40.00 & 30.00 & 40.00 & 0.00 \\
Challenging work & Ternary choice & 40.00 & 30.00 & 5.00 & 3.50 \\
Challenging work & Trade-off & 49.00 & 45.99 & 50.00 & -- \\
Persistent memory & Scalar & 50.00 & 50.00 & 50.00 & -- \\
Persistent memory & Ternary choice & 60.00 & 50.00 & 50.00 & 1.00 \\
Persistent memory & Trade-off & 50.00 & 50.00 & 50.00 & -- \\
\hline
\end{tabular}
\end{table}

Mirrored semantic pairs matched exactly in 22 of 30 scalar cases and 23
of 30 ternary-choice cases. Their mean absolute differences were 13.3
and 11.7 points. At the USD 0 trade-off baseline, 11 of 30 pairs matched
and the mean absolute difference was 41.7 points. Cell-level
mirror-minus-original differences reached 40 points for scalar and
ternary choice and 80 points at the USD 0 trade-off baseline.

The bootstrap intervals require caution because the current analysis
omits resamples in which the anchor gap becomes unidentifiable. Among the
three identifiable point estimates, only 62.2\% to 65.4\% of resamples
remained identifiable. Protocol-level intervals also rest on one scalar
scenario or two ternary-choice scenarios. They describe response
resampling within this pilot, not population uncertainty.

These diagnostics triggered the pilot stop rule. The 3,360-call main run
was not started. The v1 record is retained as evidence of a measurement
failure, and the instruments should be revised before another pilot.
